\documentclass[aspectratio=169,10pt]{beamer}

% --- Schriftarten & Sprache ---
\usepackage[utf8]{inputenc}
\usepackage[T1]{fontenc}
\usepackage[ngerman]{babel}
\usepackage{lmodern}

% --- Mathematik & TikZ ---
\usepackage{amsmath,amssymb}
\usepackage{tikz}
\usetikzlibrary{shapes.geometric,arrows.meta,positioning,calc,fit,backgrounds}
\usepackage{adjustbox}

% --- Theme-Einstellungen ---
\usetheme{Madrid}
\usecolortheme{whale}
\setbeamertemplate{navigation symbols}{}

\title{Präzisionsspektroskopie im EABC-Phasenraum}
\subtitle{Laborpfad II: Doppler-freier 2-Photonen-Polarisationskanal}
\author{Thomas Hofbauer}
\institute{Bamberg}
\date{27. September 2026}

\begin{document}

% ==============================================================================
% FOLIE: PARAMETER UND FEHLERGRENZEN VON LABORPFAD II
% ==============================================================================
\begin{frame}[t]{Laborpfad II: Spezifikation des Doppler-freien $2 \times 243$\,nm-Nulltests}
    \framesubtitle{Geometrischer Nachweis der $\mathbf{27}_{G_2}$-Tensorstruktur und Raten-Skalierung}

    \begin{center}
        % VARIANTE A: Falls als vorkompiliertes PDF vorhanden (beste Performance)
        % \includegraphics[width=0.96\textwidth,height=0.76\textheight,keepaspectratio]{laborpfad2_tikz.pdf}

        % VARIANTE B: Direktes Rendering via adjustbox (automatische Skalierung)
        \begin{adjustbox}{max width=0.98\textwidth, max height=0.77\textheight}
        \begin{tikzpicture}[
            font=\sffamily,
            >=Stealth,
            node distance=1.0cm and 0.8cm,
            every node/.style={align=center},
            header/.style={
                rectangle, rounded corners=3pt, fill=blue!80!black, text=white,
                font=\bfseries\normalsize, minimum width=4.8cm, minimum height=0.75cm
            },
            block/.style={
                rectangle, rounded corners=4pt, draw=black!60, fill=white,
                minimum width=4.8cm, minimum height=1.05cm, inner sep=5pt, font=\small
            },
            accentblock/.style={
                rectangle, rounded corners=4pt, draw=red!70!black, fill=red!5,
                minimum width=4.8cm, minimum height=1.05cm, inner sep=5pt, font=\small
            },
            mathhighlight/.style={
                font=\bfseries\small, text=blue!80!black
            },
            boundtag/.style={
                font=\bfseries\small, text=red!70!black
            }
        ]

        % SPALTE 1: LASER, OPTIK & POLARISATION
        \node[header] (col1_head) {1. Optischer Aufbau \& Geometrie};

        \node[block, below=0.5cm of col1_head] (opt1) {
            \textbf{Gegenläufige 2-Photonen-Resonanz}\\
            $2 \times 243.135\,\text{nm}$ (doppler-frei)\\
            \textsf{\footnotesize Übergang: $1S_{1/2} \to 2S_{1/2}$}
        };

        \node[block, below=of opt1] (opt2) {
            \textbf{Polarisations-Extinktion}\\
            Glan-Thompson: $> 10^{-7}$\\
            \textsf{\footnotesize Kaskaden-Filterung: $\le 10^{-12}$}
        };

        \node[block, below=of opt2] (opt3) {
            \textbf{Winkel-Konfigurationen}\\
            $\theta = 90^\circ$: Reiner Nullkanal ($\boldsymbol{\epsilon}_1 \perp \boldsymbol{\epsilon}_2$)\\
            $\theta = 45^\circ$: Kohärente Interferenz
        };

        % SPALTE 2: EXPERIMENTELLE SCHRANKEN & REGIME
        \node[header, right=1.2cm of col1_head] (col2_head) {2. Signalskalierung \& Schranken};

        \node[accentblock, below=0.5cm of col2_head] (reg1) {
            \textbf{Reiner Nullkanal ($\theta = 90^\circ$)}\\
            $\Gamma_\perp \propto \eta_A^4 \cdot \left(\frac{110}{63}\right)$\\
            \span \node[boundtag] {Nachweisgrenze: $\eta_A \lesssim 1.0 \times 10^{-2}$};
        };

        \node[accentblock, below=of reg1] (reg2) {
            \textbf{Interferenz-Kanal ($\theta = 45^\circ$)}\\
            $\Delta\Gamma \propto \eta_A^2 \cdot 2\operatorname{Re}[\mathcal{A}_{\rm SM}^* a_{27}]$\\
            \span \node[boundtag] {Nachweisgrenze: $\eta_A \lesssim 3.2 \times 10^{-3}$};
        };

        \node[block, below=of reg2] (reg3) {
            \textbf{Standardmodell-Untergrund}\\
            $\mathcal{R}_\perp^{\rm SM} \sim 2.5 \times 10^{-5} \cdot \mathcal{A}_\parallel$\\
            \textsf{\footnotesize Relativistisches $E1\text{--}M1$ \& HFS-Mischung}
        };

        % SPALTE 3: ALGEBRAISCHE FINGERABDRÜCKE (G2)
        \node[header, right=1.2cm of col2_head] (col3_head) {3. Gruppentheorie \& Signatur};

        \node[block, below=0.5cm of col3_head] (alg1) {
            \textbf{Fano-$2$-Design-Gate}\\
            $\left.\mathcal{M}_2^{\rm Haar}\right|_{\operatorname{Sym}^2(7)} = \frac{4}{7}\Pi_{\mathbf{1}} \oplus \frac{22}{63}\Pi_{\mathbf{27}}$\\
            \span \node[mathhighlight] {Universalfaktor: $c_{\mathbf{27}} = \frac{22}{63}$};
        };

        \node[block, below=of alg1] (alg2) {
            \textbf{Clebsch-Gordan-Intertwiner $\Xi$}\\
            $\mathbf{27}_{G_2} \longrightarrow \operatorname{Sym}_0^2(\mathbb{R}^3) \quad (J=2)$\\
            \span \node[mathhighlight] {Trace-Hebelarm: $\operatorname{Tr}(\Xi \mathcal{M}_2) = \frac{110}{63}$};
        };

        \node[block, below=of alg2] (alg3) {
            \textbf{Winkelcharakteristik}\\
            Signalmuster: $\beta \cdot \sin^2(2\theta)$\\
            \textsf{\footnotesize Eindeutige $d$-Wellen-Torsionssymmetrie}
        };

        % UNTERE LEISTE: SYSTEMATIK-KONTROLLE
        \node[rectangle, rounded corners=4pt, draw=black!40, fill=gray!10, 
              below=0.7cm of opt3.south west, anchor=north west, 
              minimum width=16.8cm, minimum height=1.1cm, inner sep=6pt] (sysbar) {
            \parbox{16.4cm}{\footnotesize
            \textbf{Systematische Kontrollmaßnahmen im Vakuumapparat:}\\
            $\bullet$ \textbf{Stark-Kompensation:} Aktive 3D-Helmholtz-Regelung drückt elektrostatische Streufelder auf $< 1\,\text{mV/cm}$.\\
            $\bullet$ \textbf{Laser-Stabilisierung:} PDH-Lock auf kryogenen Siliziumresonator sichert Sub-Hz-Linienbreiten.\\
            $\bullet$ \textbf{Rest-Extinktion:} Instrumenteller $10^{-7}$-Parallelanteil wird über die $\Gamma(\theta)$-Fitkurve kalibriert.}
        };

        % VERBINDUNGEN
        \draw[->, thick, color=blue!70!black] (opt1) -- (opt2);
        \draw[->, thick, color=blue!70!black] (opt2) -- (opt3);

        \draw[->, thick, color=red!70!black] (reg1) -- (reg2);
        \draw[->, thick, color=red!70!black] (reg2) -- (reg3);

        \draw[->, thick, color=blue!70!black] (alg1) -- (alg2);
        \draw[->, thick, color=blue!70!black] (alg2) -- (alg3);

        \draw[->, dashed, thick, color=black!60] (opt3.east) -- (reg1.west);
        \draw[->, dashed, thick, color=black!60] (alg2.west) -- (reg1.east);

        % HINTERGRÜNDE
        \begin{scope}[on background layer]
            \node[fill=blue!3, rounded corners=5pt, inner sep=4pt, fit=(col1_head) (opt3)] {};
            \node[fill=red!3, rounded corners=5pt, inner sep=4pt, fit=(col2_head) (reg3)] {};
            \node[fill=blue!3, rounded corners=5pt, inner sep=4pt, fit=(col3_head) (alg3)] {};
        \end{scope}

        \end{tikzpicture}
        \end{adjustbox}
    \end{center}
\end{frame}

\end{document}